\documentclass[Screen16to9,17pt]{foils}
\usepackage{zencurity-slides}

\selectlanguage{english}

\usepackage{tabularx}

% ---------------------------------------------------------------------------
% Academy 2026 - Offensive Security subject
% Day 1 (first half): Network Attacks I -- concepts
% Maps to Subject plan objective: 9 (network-attack concepts)
% Goals + Plan filled from the subject plan; content left as FIXME.
% ---------------------------------------------------------------------------

\begin{document}

\mytitlepage{Network Attacks I -- Concepts}
{Saint Cyber -- Cyber Akademiet, 2026}

\slide{Goals for today}
\begin{list2}
\item Explain common network-attack concepts: botnets and DoS/DDoS
\item Explain amplification and reflection attacks
\item Explain SYN, UDP and ICMP flooding
\item Explain ARP poisoning and man-in-the-middle risks
\end{list2}

\slide{Plan for today}
\begin{list2}
\item Attack overview and botnets
\item Denial of Service -- DoS and DDoS
\item Amplification and reflection
\item Flooding techniques across the TCP/IP layers
\item ARP poisoning and man-in-the-middle
\end{list2}

\slide{Time schedule}
% FIXME Keld: set timings for the half day
\begin{list2}
\item \textbf{FIXME:} 1) Botnets + DoS/DDoS -- XX min
\item \textbf{FIXME:} 2) Amplification / reflection -- XX min
\item \textbf{FIXME:} 3) Flooding techniques -- XX min
\item \textbf{FIXME:} 4) ARP poisoning / MITM -- XX min
\end{list2}

\slide{Types of attacks -- overview}
\begin{list2}
\item \textbf{FIXME:} where network attacks sit among the attack types
\end{list2}

\slide{Botnets}
\begin{list2}
\item \textbf{FIXME:} what a botnet is, how it is built and rented
\item \textbf{FIXME:} example (e.g. Mirai) -- keep it accurate
\end{list2}

\slide{Denial of Service (DoS)}
\begin{list2}
\item \textbf{FIXME:} definition and the shop-door analogy
\end{list2}

\slide{Distributed Denial of Service (DDoS)}
\begin{list2}
\item \textbf{FIXME:} many sources, scale, impact
\end{list2}

\slide{Amplification \& reflection}
\begin{list2}
\item \textbf{FIXME:} how reflection works (spoofed source)
\item \textbf{FIXME:} amplification example (e.g. DNS) -- verify factor
\end{list2}

\slide{TCP/IP refresher -- the layers}
\begin{list2}
\item \textbf{FIXME:} the layers relevant to the attacks below
\end{list2}

\slide{SYN flood}
\begin{list2}
\item \textbf{FIXME:} half-open connections, the three-way handshake abuse
\end{list2}

\slide{UDP flood}
\begin{list2}
\item \textbf{FIXME:} connectionless flooding
\end{list2}

\slide{ICMP flood}
\begin{list2}
\item \textbf{FIXME:} ping flood / smurf (if you cover it)
\end{list2}

\slide{ARP poisoning \& man-in-the-middle}
\begin{list2}
\item \textbf{FIXME:} how ARP poisoning redirects traffic
\item \textbf{FIXME:} the MITM risk that follows
\end{list2}

\slide{Attack landscape}
\begin{list2}
\item \textbf{FIXME:} current stats / where attacks against DK come from
\end{list2}

\slide{Summary}
\begin{list2}
\item \textbf{FIXME:} recap the concepts before the hands-on half
\end{list2}

\end{document}
